%%% Abstract.tex --- 
%% 
%% Filename: Abstract.tex
%% Author: Jacob Ringfjord & Max Wilén
%%% Code:

\begin{comment}
Applications of forensic interest are often analyzed to examine functionality that may affect the results of digital forensic investigations. When new versions of an application are released, there is a risk that assumptions about its functionality may no longer hold. Conducting a full analysis of every version of all relevant applications would entail an unreasonable amount of work. To counter this, an automated analysis would improve the quality of investigations performed while maintaining a reasonable workload.
This thesis aims to develop a method for automating the version comparison process across various executable files. The developed method is intended to help assess whether the changes in the source code are significant enough to warrant a manual analysis.

This thesis is conducted on behalf of NFC (The Swedish National Forensic Centre)—a division of the Swedish Police Authority—with the goal of developing a method for automated change detection and relevance-based filtering between two software releases.
\end{comment}

\noindent
% MOTIVATION
Applications of forensic interest are frequently analyzed to understand functionalities that may impact the outcomes of digital forensic investigations. However, when new versions of an application are released, prior assumptions about its functionality may no longer be valid. Performing a full manual analysis of every version of all relevant applications would entail an unreasonable amount of work. To counter this, an automated analysis would improve the quality of investigations performed while maintaining a reasonable workload.

% CONTRIBUTION & AIM
This thesis presents a novel approach for categorizing changes identified between two versions of an Android binary. Conducted on behalf of NFC (The Swedish National Forensic Centre)—a division of the Swedish Police Authority—the thesis presents a tool that streamlines the reverse engineering process by automatically categorizing changes found, using only the two Android binaries as input.

% METHODOLOGY
The tool is built on the Ghidriff diffing engine, which leverages Ghidra to identify changes between the two binaries. For categorization, regular expressions are used to map detected operations to a set of predefined categories, based on common symbols and function names found in Android source code. Each change is marked as \textit{critical} or \textit{non-critical} based on its forensic relevance, and changes that fall outside the predefined categories are set to \textit{uncategorized}.

% FINDINGS
The performance and accuracy of the developed tool were evaluated by comparing its output against a manual analysis of source code changes from  Signal's Android GitHub repository across three release diffs. The tool detected 91.4\% of the changes, with a categorization accuracy of 74.3\%. Across the three release comparisons used as test data, only 3.3\% of unique changes were marked as \textit{uncategorized}, indicating they did not match any predefined category. As a result, the thesis presents convincing evidence that the proposed tool is effective both in filtering out insignificant changes and in highlighting meaningful ones.


%%%%%%%%%%%%%%%%%%%%%%%%%%%%%%%%%%%%%%%%%%%%%%%%%%%%%%%%%%%%%%%%%%%%%%
%%% Abstract.tex ends here
